Nous ne citerons pas tous les axes d'amélioration car il en existe beaucoup, mais voici
une liste non exhaustive de certains points à améliorer.
\bigskip

Tout d'abord, il faut savoir que celui-ci n'est pas fini. Il manque
encore certaines fonctionnalités telles que la mise en œuvre d'une zone de sécurité
dynamique en fonction de la vitesse du véhicule.
L'accéléromètre a été implanté sur la carte, mais le code non.

Le coût du prototype peut également être réduit de manière conséquente. À l'heure actuelle,
le projet atteint les 600 euros, mais cela s'explique par des fournisseurs imposés.
Le prix de certains composants est donc monté à 75 euros au lieu de 35 euros (le tout x5).

Les délais étant très serrés, il manque également de nombreuses fonctionnalités, tant sur le principe
de détection des ouvriers que sur l'application web. 

L'impression PCB est également à améliorer, pour passer de cartes prototypes à de vraies cartes industrielles.
Cela permettrait d'éviter de potentiels court-circuits grâce au vernis, les problèmes de soudure et d'offrir un gain de temps.
Nous aurions également une facilité de debug grâce à la sérigraphie.

Un autre point d'amélioration serait la mise en œuvre de plots anti-vibrations sur les cartes.
Cela permettrait de réduire les chocs sur les cartes lorsque le véhicule est en activité.

Nous avons également pensé à la mise en place de gaines thermo-rétractables au niveau des bus CAN,
pour assurer une étanchéité de nos câbles.

Il est également envisageable de revoir certains composants, tels que l'abaisseur XL0415 \cite{XL0415}, qui 
se trouve être surdimensionné pour la carte et les besoins de celle-ci.

Un autre point important est la mise en place des ancres. Dans son fonctionnement idéal,
les ancres arrivent à connaitres leurs positions par rapport au plan entre elles.
Mais à l'heure actuelle, l'emplacement des ancres se fais en hard, n'ayant pas eu le temps
de mettre en place le code. 

